\section{Significance of the Study}
Hospitals are essential facilities whose seismic performance has consequences beyond the direct repair of structural damage. By comparing fixed-base and flexible-base representations within each of the six structural-block datasets, this study provides a consistent basis for examining changes in modal characteristics under differing foundation assumptions in the Nepalese context.

The analysis considers fundamental period, modal participation, displacement, drift, base and storey shear, foundation deformation, member forces and fragility relationships. Considering several response measures is useful because period changes alone do not determine all displacement, force or damage outcomes.

\section{Scope and Limitations of the Study}
The study considers the six G+4 structural block models with fixed, hard-, medium- and soft-soil foundation supports. Modal and response-spectrum results are investigated for seismic demand, whereas foundation displacement and structural member force envelopes are investigated for the specified gravity-load combination. The fragility assessment is preliminary and applies to the damage-state definitions used for the analytical curves.

\section{Organization of the Study}
The thesis is organized into five chapters. Chapter 1 defines the research problem, questions, objectives, significance, scope, and limitations. Chapter 2 reviews literature on seismic vulnerability, hospital performance, SSI, performance-based assessment, and fragility. Chapter 3 describes the structural models, foundation-support conditions and methods of analysis. Chapter 4 presents and discusses the seismic and foundation response of the six blocks. Chapter 5 summarizes the principal conclusions, recommendations, and limitations of the study.
